%% ma: L12-B07-0028
%% lop: 12
%% chuong: 2
%% bai: 7
%% dang: 12.7.10
%% loai: TN4
%% muc: TH
%% dap_an: "A"
%% nguon: "CĐ2.2.pdf giáo viên gửi 27/10/2026 (M07, Lê Minh Khôi – Nguyễn Anh Khôi, Chương II Chủ đề 2 Tọa độ của vectơ trong không gian), Câu 30"
%% kien_thuc: "Bài 7 (tọa độ điểm khi gắn hệ trục), lớp 11 (hai mặt phẳng vuông góc)"
%% trang_thai: cho-duyet
%% ly_do: "Dạng mới 12.7.10 đề xuất, chờ giáo viên duyệt"
%% ghi_chu: "Bổ sung câu chữ cách chọn trục theo hình gốc"
\begin{ex}
Cho hình chóp $S.ABCD$ có đáy $ABCD$ là hình vuông cạnh bằng $1$, $SAD$ là tam giác đều và nằm trong mặt phẳng vuông góc với đáy. Gọi $H$ là trung điểm $AD$. Thiết lập hệ tọa độ $Oxyz$ có gốc tọa độ trùng với điểm $H$, tia $Ox$ chứa $D$, tia $Oy$ cùng hướng với $\overrightarrow{AB}$, tia $Oz$ chứa $S$ (như hình vẽ).
\begin{center}
\begin{tikzpicture}[scale=2.4,x={(-.55cm,-.4cm)},y={(1.3cm,0cm)},z={(0cm,1cm)}]
  \path (0,0,0) coordinate (H) (-.5,0,0) coordinate (A) (.5,0,0) coordinate (D) (.5,1,0) coordinate (C) (-.5,1,0) coordinate (B) (0,0,.866) coordinate (S);
  \draw (S)--(D)--(C)--(B)--cycle (S)--(C);
  \draw[khuat] (S)--(A) (D)--(A)--(B) (S)--(H);
  \draw[truc,red] (D)--++(.5,0,0) node[left]{$x$};
  \draw[truc,red] (H)++(0,1,0)--++(0,.5,0) node[right]{$y$};
  \draw[khuat] (H)--++(0,1,0);
  \draw[truc,red] (S)--(0,0,1.3) node[right]{$z$};
  \foreach \P in {H,A,B,C,D,S} \fill (\P) circle (.5pt);
  \path (S) node[above left]{$S$} (A) node[above]{$A$} (B) node[right]{$B$} (C) node[below]{$C$} (D) node[below]{$D$} (H) node[above left]{$H$};
\end{tikzpicture}
\end{center}
Biết rằng tọa độ các điểm $S$, $C$ là $S(x_S;y_S;z_S)$, $C(x_C;y_C;z_C)$. Tính giá trị $L=x_C+y_C-z_S^2$.
\choice{\True $L=0{,}75$}{$L=0{,}5$}{$L=1$}{$L=0$}
\loigiai{$SH\perp(ABCD)$, $SH=\dfrac{\sqrt3}2$. $C\left(\dfrac12;1;0\right)$, $S\left(0;0;\dfrac{\sqrt3}2\right)$.

$L=\dfrac12+1-\dfrac34=0{,}75$.}
\end{ex}
